%=============================================================================!
% 4.8  TWO BODIES AND THE REDUCED MASS                                        !
%=============================================================================!
\section{Two bodies and the reduced mass}
\label{sec:os-bodies}

So far one end of every spring has been fixed to a wall. Two atoms joined
by a bond, or the two carts of Chapter~2, have no wall: both ends move.
This section shows that their motion splits into a steady drift of the
pair and an oscillation of the distance between them, with a mass smaller
than either.

\subsection*{The relative motion}

Two bodies of masses $m_1$ and $m_2$, at positions $x_1 < x_2$ on a line,
are joined by a spring of stiffness $k$ and natural length $\ell$. Its
stretch is $x_2 - x_1 - \ell$, so it pulls the first body forwards and the
second backwards with forces of size $k(x_2 - x_1 - \ell)$:
\begin{equation*}
   m_1\ddot{x}_1 = +k\,(x_2 - x_1 - \ell) , \qquad
   m_2\ddot{x}_2 = -k\,(x_2 - x_1 - \ell) .
\end{equation*}
No external force acts along the line, so by \cref{sec:ne-bodies} the
centre of mass moves at a constant velocity. For the separation $r = x_2 -
x_1$, divide each equation by its mass and subtract the first from the
second:
\begin{align*}
   \ddot{r} = \ddot{x}_2 - \ddot{x}_1
   &= -\frac{k}{m_2}\,(r - \ell) - \frac{k}{m_1}\,(r - \ell)
   = -k\Bigl(\frac{1}{m_1} + \frac{1}{m_2}\Bigr)(r - \ell) .
\end{align*}
Defining the reduced mass $m_{\mathrm r}$ by $1/m_{\mathrm r} = 1/m_1 + 1/m_2$, or, putting the
fractions over a common denominator and inverting,
\begin{equation}
   m_{\mathrm r} = \frac{m_1 m_2}{m_1 + m_2} ,
   \label{eq:os-reduced}
\end{equation}
the separation obeys $m_{\mathrm r}\ddot{r} = -k(r - \ell)$. Writing $s =
r - \ell$ for the stretch, $\ddot{s} = \ddot{r}$ since $\ell$ is constant,
so $m_{\mathrm r}\ddot{s} = -ks$: the stretch oscillates about zero like a
single body of mass $m_{\mathrm r}$ on a spring fixed to a wall
(\cref{sec:ne-spring}), with $\omega = \sqrt{k/m_{\mathrm r}}$.

The carts of \cref{sec:ne-bodies}, of \SI{1}{\kilogram} and
\SI{3}{\kilogram}, have $m_{\mathrm r} = 3/4 = \SI{0.75}{\kilogram}$, and with the
spring of \SI{200}{\newton\per\metre} the separation would oscillate with
$\omega = \sqrt{200/0.75} = \SI{16.33}{\radian\per\second}$. Released at
rest with the spring compressed, the separation starts at a turning point
and reaches the natural length a quarter of a period later, at
$\frac14\times2\pi/\omega = \SI{0.0962}{\second}$. This is the moment,
found in \cref{sec:ne-bodies} by computer, at which the spring fell away.

\subsection*{Light partners do the moving}

The reduced mass is smaller than either mass, since $m_{\mathrm r} =
m_1\times m_2/(m_1 + m_2)$ with $m_2/(m_1 + m_2) < 1$, and likewise
$m_{\mathrm r} = m_2\times m_1/(m_1 + m_2)$ with $m_1/(m_1 + m_2) < 1$. When one body is much heavier, it is close to
the lighter mass: for an oxygen atom and a hydrogen atom,
\begin{equation*}
   m_{\mathrm r} = \frac{1.008\times15.999}{1.008 + 15.999} = \SI{0.9483}{\amu} ,
\end{equation*}
\SI{94}{\percent} of the hydrogen mass. For a molecule that is not drifting,
the centre of mass stays still, so $m_{\mathrm H}x_{\mathrm H} + m_{\mathrm
O}x_{\mathrm O}$ keeps its value and any changes obey $m_{\mathrm
H}\,\Delta x_{\mathrm H} = -m_{\mathrm O}\,\Delta x_{\mathrm O}$: the
hydrogen moves $15.999/1.008 = 15.9$ times as far as the oxygen, the other
way: the vibration
of an O--H bond is mostly the hydrogen atom moving against an almost still
partner. For the same spring, a smaller mass means a higher frequency,
$\omega = \sqrt{k/m_{\mathrm r}}$, which is why bonds to hydrogen vibrate fastest
(\cref{sec:os-molecules}). For two atoms in a Morse or Lennard-Jones well
(\cref{sec:os-anharmonic}), the same argument applies with $r$ the
distance between them and $m_{\mathrm r}$ in place of $m$, for motion
along their line of centres without rotation.

\begin{selfcheck}
Two equal masses $m$ are joined by a spring of stiffness $k$. What is the
reduced mass, and with what angular frequency does the separation
oscillate?
\end{selfcheck}
